\documentclass{article}
\usepackage[T1]{fontenc}
\usepackage{lmodern}
\usepackage{graphicx}
\usepackage{amsmath,amssymb}
\usepackage[a4paper,margin=2cm]{geometry}
\usepackage[absolute,overlay]{textpos}
\setlength{\TPHorizModule}{1mm}
\setlength{\TPVertModule}{1mm}
\usepackage[indonesian]{babel}
\usepackage{hyperref}
\setlength{\emergencystretch}{2em}
% unit: Halaman 21

\begin{document}

% === Halaman 21 (llm) ===

\section*{Jelaslah Kurva Eliptik membentuk Grup $\langle G, + \rangle$}

Karena:
\begin{itemize}
    \item Himpunan $G$: semua titik $P(x,y)$ pada kurva eliptik
    \item Operasi biner: $+$
    \item Semua aksioma terpenuhi sbb:
    \begin{enumerate}
        \item Closure: semua operasi $P + Q$ berada di dalam $G$
        \item Asosiatif: $P + (Q + R) = (P + Q) + R$
        \item Elemen netral adalah $O$: $P + O = O + P = P$
        \item Elemen invers adalah $-P$: $P + (-P) = O$
        \item Komutatif: $P + Q = Q + P$ \textcolor{red}{(abelian)}
    \end{enumerate}
\end{itemize}

\end{document}
